\documentclass[11pt,a4paper]{article}

% --- Packages ---
\usepackage[T1]{fontenc}
\usepackage[english]{babel}      % use [italian] for Italian text
\usepackage{amsmath, amssymb, amsthm, mathtools}
\usepackage{geometry}
\usepackage{booktabs}
\usepackage{eurosym}
\geometry{margin=2.5cm}


\title{InvoiceFlow\\[0.5em] \large Software Engineering project}
\author{Lorenzo Arduini and Riccardo Striano}
\date{\today}

\begin{document}
\maketitle

\section*{What the project does}

InvoiceFlow reads the invoices a freelancer or a small business receives, checks that each one
is correct, finds the ones uploaded more than once and flags suspicious ones, such as a known
supplier suddenly asking to be paid to a new bank account. It reads Italian e-invoices
(FatturaPA XML, including signed \texttt{.xml.p7m} files) and Italian PDF invoices. It is
written in Python and runs locally: from the command line, or through a Streamlit web app on
\texttt{localhost} that shows spending, payments coming due and every problem found, each with
an explanation. Running locally was deliberate: invoices contain names, VAT numbers, IBANs and
addresses, so nothing leaves the user's machine, and we tested on real invoices only after
anonymizing them.

\section*{General components}

A pipeline: importers turn each file into one shared invoice model, the checks run on it at
import time, and the result goes into a SQLite archive, which duplicate and fraud detection read
as a whole and the dashboard displays. The checks are eight rules from the FatturaPA specification,
verifying that the amounts add up (lines, VAT, total, payments), that VAT numbers and IBANs have
valid check digits and that due dates make sense. Fraud detection has three rules, each comparing
an invoice with its supplier's history: a new IBAN, a look-alike supplier name, an unusually
large amount. The README lists every rule exactly. Beside the pipeline sit an anonymizer and an invoice generator, which builds a dataset
of fake invoices from a seed, with planted errors, frauds and duplicates and an answer key
(\texttt{labels.csv}), so every run can be measured and reproduced.

\section*{Hurdles and interesting problems}

\textbf{No two invoices fill the same fields.} FatturaPA has hundreds of fields, mostly optional,
and every invoicing program fills a different subset, so a parser written around one supplier
fails on the next. We first agreed on a small shared model holding only what the checks and
fraud rules need, and wrote every reader (XML, \texttt{.p7m}, PDF) to fill it, treating a
missing optional field as normal rather than as an error.

\textbf{Correct invoices that look wrong.} A field the program does not understand can make a
legitimate invoice look broken. An invoice with 0\% VAT above \euro{}77.47 carries a \euro{}2
stamp duty (\emph{bollo}) added to the total, so the total no longer equals taxable amount plus
VAT; withholding tax did the same to payments. Each time we had to choose between loosening the
rule and teaching the model the new field. We kept the arithmetic exact to the cent and added
the most common fields; stamp duty and discounts are still a known source of warnings.

\textbf{Rules a fraudster can exploit.} With seed 2026 the fraud rules raised 475 false alarms
on 9,498 regular invoices (5 per 100, against 0.01 with other seeds), 472 of them on one genuine
supplier, Sudsoftware S.a.s. The look-alike rule takes the supplier it sees first as genuine,
and the fake ``uSdsoftware S.a.s.'' had an invoice on the same day as the real one's first,
imported earlier. Everything was inverted: the fake became the trusted supplier and every real
invoice was reported as an impersonation. A fraudster who gets one invoice into the archive
first gets the same result; judging who is genuine by a longer history would make this much
harder. The unusual-amount rule flags totals above 2.5 times the supplier's largest earlier
invoice, but every invoice that passes raises that bar: from a usual \euro{}1,000, invoices of
\euro{}2,500, \euro{}6,250, \euro{}15,625 and \euro{}39,000 all pass unnoticed.

\section*{Ideas}

\textbf{Testing against invoices we made ourselves.} Real invoices with known frauds do not
exist, so we measured the program against the generator's answer key. We tuned the rules on seed 42 and
measured on a fresh seed, 12345: 10,000 invoices, with 464 errors, 472 frauds and 300
duplicates. The checks caught all 464 errors, of all eight kinds, with no false alarm on 9,536
invoices that should pass. Of the 342 frauds in scope, 317 were flagged (92.7\%): all look-alike
suppliers and IBAN changes, and 98 of 123 unusual amounts, with one false alarm on 9,528 regular
invoices. The other 130 use non-existent VAT numbers, which only the EU VIES service can detect.

\textbf{Anonymizing without breaking the rules.} Random replacements would have made real
invoices useless for testing, since the fraud rules depend on recognizing a supplier and its
usual IBAN. With the same secret key, a real supplier always gets the same fake name, VAT number
and IBAN, with check digits valid exactly when the original's were, so every check gives the
same answer on the copy.

\section*{Looking ahead}

With a more precise invoice model, one that understands and can verify every field, InvoiceFlow
could become a tool used by an AI agent set up to help a small business with its finances:
reading every incoming invoice, warning before a suspicious one is paid, and keeping track of
what is coming due.

\end{document}
